\documentclass{report}
\usepackage{amsmath,amssymb}
\setlength{\parindent}{0mm}
\setlength{\parskip}{1em}
\begin{document}
\begin{center}
\rule{6in}{1pt} \
{\large Raya Horesh \\
{\bf Heat Transfer Model and Inversion for Smarter Buildings}}

IBM T J Watson Research Center \\ 1101 Kitchawan Rd \\ Yorktown Heights \\ NY 10598
\\
{\tt rhoresh@us.ibm.com}\\
Lianjun An\\
Young M. Lee\\
Young T. Chae\\
Rui Zhang\end{center}

In the United States, the energy consumption of commercial and
residential buildings exceeds 40% of the total demand, while being
responsible for 45% of the total greenhouse gas emissions. As a
consequence, saving related energy and costs, improving energy
consumption efficiency and reducing greenhouse emissions are becoming key
initiatives in many cities and municipalities. For that reason it also
plays an important role in IBM smarter planet initiative.\\
In order to reduce energy consumption in buildings, one needs to
understand and be capable of modeling the underlying heat transfer
mechanisms, characteristics of building structures, operations and
occupant energy consumption behaviors. \\
Often, some of the crucial building envelope parameters are not known. In
order to infer thermal parameters associated with building�s envelope,
sensor data is utilized within an inversion procedure. The forward
problem involves a system of ODEs capturing the governing heat transfer
model. Our formulation also involves derivation of adjoints for efficient
gradient computation. Through sensitivity analysis, we assess the impact
of potential energy saving retrofits and their quantitative impact upon
energy consumption of commercial buildings.


\end{document}
